\section{Introduction}\label{sec:introduction}

A private-spend protocol can be mathematically specified, machine-checked,
implemented, compiled, and deployed while still leaving gaps between those
stages.  A kernel theorem does not identify the source that was compiled; a
source proof does not identify the deployed bytes; and a finalized transaction
does not by itself establish either mathematical security or source
correspondence.  The central contribution of Aspis is therefore an evidence
chain, not an isolated benchmark: maintained private-spend mathematics in
Lean~4, selected production Rust connected to that mathematics, a
byte-reproducible Solana binary, and a finalized mainnet state transition.

The chain is intentionally explicit about direction and scope.  Lean
kernel-checks the relation, finite security arithmetic, circle and masking
lemmas, and their composition \cite{deMouraUllrich2021Lean4}.  For the selected
    V5 release schedule, Charon extracts selected production Rust and Aeneas
    translates it into functional Lean \cite{HoProtzenko2022Aeneas}; bridge
    proofs cover Component~A at the selected schedule, the Component-B/C
    evaluator and packed-output paths, and Tag-67 instruction decoding plus the
    six ordered work checks.  The result is not a universal refinement theorem
    for every schedule or verifier path.  Its one
remaining Tag-67 implementation boundary equates the extracted
hash-function application with the maintained transcript hash on
\(\texttt{DOM\_GRIND}\mathbin\Vert\texttt{nonce\_le64}\).  A clean checkout of
the identified source reproduces the 1,258,496-byte V5 SBF binary
byte-for-byte.  The same binary was deployed on mainnet-beta, where exact
signed-wire simulation and landed execution of a 75,358-byte proof both
consumed 1,334,452 compute units.

Transparent arguments provide the cryptographic setting for this chain.  They
remove the structured reference string and secret trapdoor required by many
succinct proof systems.  FRI and subsequent STARK constructions showed how
low-degree testing and hash-based commitments can support transparent
verification of large computations
\cite{BenSassonBentovHoreshRiabzev2018FRI,
BenSassonBentovHoreshRiabzev2018STARK}.  Deploying such a verifier inside a
blockchain transaction adds a distinct constraint: it must fit the runtime's
transaction budget while preserving application state-transition semantics.
Solana meters execution in compute units and stores persistent state in
explicitly addressed accounts \cite{SolanaComputeBudget,SolanaAccounts}.  A
proof that verifies cheaply on a workstation, or even a source-level verifier
that has been proved correct, is not by itself a deployed on-chain
construction.

Yano's measurement study \cite{Yano2025Measurement} showed that the budget is
not prohibitive: a minimal Winterfell STARK verifies inside one Solana devnet
transaction.  Aspis began by reproducing that measurement and modelling its
costs, then asked whether the same budget could hold a complete shielded-spend
statement and atomic state transition at deployable parameters on
mainnet-beta.  The manuscript answers that question with a frozen q18/g37
security case study; the later V5 Tag-67 release adds the selected-source
correspondence and byte-reproducible deployment result.  The two are distinct
artefacts.  In particular, the q18/g37 numerical security bounds in this paper
are not silently inherited by V5.

Both releases study one deliberately narrow application atom.  A valid spend
replaces one note by one output note at the same private path in a depth-20
tree, preserves the asset identifier, enforces value conservation including a
public fee, derives a nullifier, and advances a public pool anchor and sequence
number.  The program accepts only if the proof verifies against that public
state; in the same transaction it creates the nullifier marker and changes the
pool state.  This follows the commitment-and-nullifier structure used by
shielded payment systems \cite{BenSassonEtAl2014Zerocash}, but it is not a
general wallet protocol, multi-input transaction, deposit path, or public
append operation.

The central systems choice is to separate proof transport from proof use.  A
proof is first written to a program-owned account and finalized.  A later
instruction reads the finalized bytes, verifies the complete proof, and
performs the state mutation atomically.  Thus, ``one transaction'' means the
verification and state update from a finalized, pre-uploaded proof account; it
does not include account creation, upload, or finalization.  The separation
keeps proof bytes available to the verifier without implying that the full
proof lifecycle fits in one transaction.

\paragraph{Protocol.}
\Profile{} uses arithmetic over $\Fbase$ and a degree-four extension field.
Its commitment layer is a custom, circle-domain, WHIR-style multilinear
protocol.  The term \emph{WHIR-style} identifies the use of batched
Reed--Solomon proximity testing, out-of-domain checks, recursive folding, and
query openings in the lineage of WHIR
\cite{ArnonChiesaFenziYogev2025WHIR}; it does not assert
protocol equivalence.  The frozen instance has rate $1/512$,
$q=\ProfileQueryCount$ sampled query fibers, a batch-work parameter
$g=\ProfileBatchWork$, two commitment layers, and four arity-four folds.  A
three-branch selector chooses the least schedule satisfying the exact
\(\mathsf{Good}_{\mathrm{spend}}\) rank conditions.  The transcript, selector rule,
serialization,
and boundary count are part of the protocol definition rather than prover
heuristics.

\paragraph{Security.}
The soundness argument is stated in the random-oracle model (ROM) and begins
with an explicit interactive event ledger.  It
specializes proven Reed--Solomon proximity and mutual-correlated-agreement
results to the circle code, batching map, folds, out-of-domain checks, and
query schedule used by the implementation.  It then applies a 32-boundary
Fiat--Shamir accounting rule for a classical random-oracle adversary.  For
$1\leq T\leq 2^{128}$ random-oracle queries, the released metric is false
acceptance probability divided by $T$; after the conservative whole-ledger
three-selector union, its minimum floor is \ProfileSoundnessBits{} bits.  The
claim is argument soundness for the exact relation.  No extractor or knowledge
soundness claim is made.

The hiding statement is likewise tied to a precise experiment.  It compares
the complete declared verifier-visible Proof-or-Abort view with an explicit
simulator in a programmable SHA-256 random oracle.  The view includes the
variable-length proof bytes, opened records, transcript and work values,
selector or Abort result, serialization, program logs, and deterministic
public mutation.  It excludes network metadata, payer linkage, timing, local
worker state, and physical side channels.  With at most $2^{128}$ adaptive
oracle queries, including post-output queries, and at most 17 proving
attempts, the real-versus-simulator floor is \ProfileSimulationBits{} bits.
The triangle inequality gives a \ProfilePairwiseBits-bit bound between the
views of any two valid witnesses for the same public statement.  This is a
computational hiding result in the stated model, not statistical honest-verifier
zero knowledge or a standard-model theorem.

\paragraph{Implementation and evidence.}
The verifier is implemented as a Solana SBF program.  Its proof account has an
application-level state machine: bytes may be uploaded before finalization,
whereas the spend instruction requires the finalized state.  Verification
precedes every state write, and the pool pre-state, sequence number, anchor,
nullifier address, account ownership, and proof finalization state are checked
at the instruction boundary.  The frozen q18/g37 release binds a
\ProfileProofBytes-byte proof to a \ProfileSbfBytes-byte SBF binary.  The
maximum measured local path consumes \ProfileComputeUnits{} compute units.  A
finalized mainnet-beta execution of that exact release consumed
\ProfileMainnetComputeUnits{} compute units, advanced the pool sequence from
zero to one, created the canonical nullifier marker, and refunded the complete
proof-account balance in the same transaction.  Deployment, account creation,
proof transport, and proof finalization occurred beforehand and are reported
separately.  The later V5 release binds the selected source-correspondence
theorem, a byte-reproducible 1,258,496-byte SBF binary, and a 75,358-byte
proof.  Its atomic Tag-67 transaction finalized at 1,334,452 compute units in
both exact signed-wire simulation and landed metadata, with canonical
nullifier-PDA bump 255.  V5 deployment, its 84-transaction proof lifecycle,
and its cleanup transactions are recorded separately from the atomic spend
and from the q18/g37 release.

\paragraph{Trust boundary.}
The Lean kernel checks the theorems presented to it, including the selected
Component-A schedule, Component-B/C evaluator and output, and Tag-67 work-wire
composition described above.  The remaining trusted or assumed boundary
includes the single transcript-hash-call equality stated above; the
Charon/Aeneas translation and bridge inputs; Rust outside those selected paths;
the Rust, LLVM, and Solana SBF compilation path; the Solana loader, account,
rollback, and compute metering semantics; SHA-256 and the random-oracle
idealisation; and the cited
proximity, agreement, and cryptographic primitive results.  Byte reproducibility
detects build divergence but is not a verified compiler theorem.  Neither the
protocol nor the implementation has received an external security audit.

\paragraph{Contributions.}
The paper makes the following contributions.

\begin{enumerate}[leftmargin=*,itemsep=0.45em]
  \item It specifies the exact relation $\RProfile$, public account pre-state,
  deterministic post-state, proof-account state machine, and byte-level
  transcript consumed by the program.

  \item It gives a finite-parameter soundness derivation for the
  $q=\ProfileQueryCount$, rate-$1/512$ circle-domain instance, including the
  event union, selector loss, work schedule, and 32-boundary random-oracle
  reduction.

  \item It constructs a witness-free simulator for the complete declared
  Proof-or-Abort view and states both the real-versus-simulator and pairwise
  witness-hiding games with their concrete bounds.

  \item It supplies a Lean~4 formalisation of the relation core, finite
  security arithmetic, circle and masking lemmas, and theorem composition,
  then connects selected V5 verifier paths through Charon/Aeneas and explicit
  bridge proofs: Component~A at the release schedule, the Component-B/C
  evaluator and public output, and the Tag-67 work-wire and six ordered checks.
  The proof-status ledger separates kernel theorems, the one
  transcript-hash-call premise, and cited cryptographic interfaces.

  \item It binds that source-linked verifier to a byte-reproducible SBF
  implementation and reports the finalized V5 mainnet-beta Tag-67 execution
  of a 75,358-byte proof at 1,334,452 compute units, while preserving the
  distinct q18/g37 execution and its quantitative security claims.

  \item It provides machine-readable ledgers, rank certificates, known-answer
  tests, proof and binary hashes, negative tests, and transaction evidence from
  which the quantitative claims can be regenerated.
\end{enumerate}

\paragraph{Organization.}
\Cref{sec:background} introduces the coding, transcript, and Solana concepts
used in the construction.  \Cref{sec:relation} defines the statement, relation,
and account model.  \Cref{sec:construction} specifies the protocol and wire
format, and \Cref{sec:formalization} records the machine-checked development.
\Cref{sec:soundness,sec:hiding} prove soundness and computational hiding.
\Cref{sec:implementation} gives the implementation refinement and atomicity
argument, and \Cref{sec:evaluation} reports the measurements and reproduction
procedure.  \Cref{sec:limitations} states the scope of the result.
